
Current AI alignment evaluation focuses on AI-to-human alignment while largely ignoring human-to-AI alignment: whether AI systems preserve human agency and critical evaluation capability. This work investigates whether existing preference data contains extractable signals related to agency preservation by introducing the Bidirectional Alignment Index (BAI), a composite metric computed from four linguistic proxies: clarifying questions, option enumeration, epistemic hedging, and explicit deferral. Experiments on HH-RLHF and RewardBench data (N=41,896) yield mixed results. Agency proxies achieve high extraction reliability (mean AUROC 0.98), and BAI survives adversarial probing designed to remove reward-predictive variance (AUROC 0.99 after gradient reversal on synthetic hidden states, reward $R^2$ change -0.27\%). However, semantic analysis reveals that high-BAI responses cluster by generic conversational patterns rather than interpretable agency vocabulary (0\% agency pattern rate across 3 discovered topics). These findings contribute both a validated methodology for testing representational independence between alignment dimensions and an important negative result: surface-level linguistic proxies, while discriminatively reliable, do not capture semantically coherent agency content. Future operationalizations of bidirectional alignment require embedding-based or LLM-as-judge approaches rather than pattern-matching heuristics.
